\section*{Chapter \ref{ch:g5:people-and-nature} --- People and Nature}

\begin{solution}{exo:g5:people-and-nature:1}
Fields, managed for one chosen crop; \omterm{prop:g5:people-and-nature:managed}{forests}, managed for
timber; meadows, kept open for hay or grazing; towns, roads
and dams, built for living and moving. (Any four with their
purposes.)
\end{solution}

\begin{solution}{exo:g5:people-and-nature:2}
Left unmowed and ungrazed, it grows back to scrub within a
few years, and then to woodland --- the open, flowery meadow
\omterm{def:g2:where-animals-live:habitat}{habitat} disappears.
\end{solution}

\begin{solution}{exo:g5:people-and-nature:3}
Take no faster than it renews; keep the workers (the \omterm{def:g5:biodiversity:species}{species}
whose free services the use depends on); return what you
borrow (close the loops, leave no piling wastes).
\end{solution}

\begin{solution}{exo:g5:people-and-nature:4}
Hedges grubbed out versus kept: A loses the soil-holders and
the homes of pest-eaters. Bare winter soil versus covered:
A's soil washes downhill in the rains. One giant single crop
versus rotation: A invites any pest of that crop to sweep the
whole hillside, and never rests the ground.
\end{solution}

\begin{solution}{exo:g5:people-and-nature:5}
Count what the forest grows in a year; cut no more than
that. It is rule 1 --- take no faster than it renews --- and
the forest yields timber every year, indefinitely, while
remaining a living forest.
\end{solution}

\begin{solution}{exo:g5:people-and-nature:6}
For example: swifts, nesting on concrete ``cliffs''; foxes,
patrolling bins and gardens; wild \omterm{def:g1:plants-around-us:plant}{plants}, rooting in asphalt
seams and wall \omterm{prop:g3:bones-and-muscles:joints}{joints} (also pigeons on ledges, falcons on
towers).
\end{solution}

\begin{solution}{exo:g5:people-and-nature:7}
For people: slower water and flood storage in the restored
wet meadows --- less flooding downstream. For the web: the
bends' \omterm{def:g2:where-animals-live:habitat}{habitats} came back, and with them \omterm{prop:g3:sorting-living-things:groups}{fish}, dragonflies
and the storks' hunting meadows.
\end{solution}

\begin{solution}{exo:g5:people-and-nature:8}
Open answer. Typical finds: mowed verges, pruned trees,
planted beds, weeded pavements --- with perhaps one
unmanaged wall-top or ditch, which is often the wildest
spot in view.
\end{solution}

\begin{solution}{exo:g5:people-and-nature:9}
Never-touching would actually destroy some \omterm{def:g2:where-animals-live:habitat}{habitats}: the
flowery meadow exists \emph{because} of mowing, and grows
shut without it. Protecting nature is not zero touch but
good management --- the honest question is how to touch:
which rules the use respects, not whether hands were laid
on at all.
\end{solution}

\begin{solution}{exo:g5:people-and-nature:10}
Rule 2 broken outright (workers evicted with the hedges),
rule 3 partly (soil washed away is never returned), and
rule 1 endangered (soil lost faster than soil forms). Most
improving single change: replant the hedges --- they hold
the soil and rehouse the pest-eaters at one stroke.
\end{solution}

\begin{solution}{exo:g5:people-and-nature:11}
The doubled catch took \omterm{prop:g3:sorting-living-things:groups}{fish} far faster than the shoals bred
--- spending the breeding stock itself. Spending capital
always looks splendid while it lasts: the holds were full
precisely because next year's \omterm{prop:g3:sorting-living-things:groups}{fish} were being landed this
year. The forester harvests only the year's growth and
keeps the growing stock whole --- same arithmetic, opposite
sign, which is why the forest yields forever and the fleet
stopped in year six.
\end{solution}
